\documentclass[11pt,a4paper]{article}
\usepackage[margin=1in]{geometry}
\usepackage{hyperref}
\usepackage{enumitem}
\usepackage{graphicx}
\usepackage{xcolor}
\usepackage{titling}

\newcommand{\BvrHint}[1]{%
  {\normalfont\normalsize\itshape\hfill #1}}

\pretitle{\begin{center}\bfseries\fontsize{18bp}{18bp}\selectfont}
\posttitle{\par\end{center}\vskip 0.5em}

\title{AI Response Analysis and Prompt Engineering Platform}

\author{Student Project Proposal}

\date{\today}

\begin{document}
\maketitle

\section{Higher-order goal%
  \BvrHint{\ldots secondary framing}}

This project aims to help users understand, evaluate, and improve the AI
responses they receive from generative AI systems. The proposed software will
primarily function as an AI response analysis and explanation platform, with
prompt engineering used as a supporting mechanism for improving unsatisfactory
responses.

The immediate engineering objective is to build a usable dashboard that
examines the quality, structure, relevance, clarity, completeness, and
constraint alignment of AI-generated responses. It will also provide
user-facing explanations connecting observable response characteristics with
the user's prompt and requested requirements. The project will not attempt to
reveal hidden model reasoning or internal chain-of-thought.

\section{Time-to-value%
  \BvrHint{\ldots primary heuristic}}

\begin{itemize}[leftmargin=0.7cm]
    \item The first iteration will provide AI response generation and a
    basic response quality analysis.
    \item Subsequent iterations will add response explanation, response
    diagnosis, response comparison, prompt analysis, and prompt improvement.
    \item Development will prioritize simple, demonstrable features that can
    be tested independently and integrated into the final dashboard.
    \item Success is defined by helping users understand why an AI response
    has its observable characteristics, identify problems, and improve the
    result when necessary.
\end{itemize}

\section{Problem Statement}

Users frequently receive AI responses that are too long, too short, too
technical, incomplete, poorly structured, irrelevant, or not aligned with
their requested format. Although users can see the final response, they may
not understand what caused these observable characteristics or how to decide
whether the response is actually suitable for their goal.

Common pain points include:

\begin{itemize}[leftmargin=0.7cm]
    \item \textbf{Limited response understanding:} users may receive an
    answer without knowing whether it is relevant, complete, clear,
    concise, well-structured, or aligned with their request.
    \item \textbf{Difficulty diagnosing poor responses:} users may know that
    an answer is unsatisfactory but may not know whether the problem comes
    from missing context, excessive detail, weak structure, or unmet
    constraints.
    \item \textbf{Lack of response transparency:} users generally see the
    final answer but do not receive a simple analysis of the observable
    factors associated with that answer.
    \item \textbf{Repeated trial and error:} users may repeatedly regenerate
    answers instead of understanding what should be changed.
    \item \textbf{Inefficient AI usage:} users with limited messages,
    responses, or usage quotas can benefit from analyzing an answer and
    improving the next request before using another interaction.
\end{itemize}

\textbf{Why this matters now (contextual relevance):} Generative AI is being
used for education, programming, research, writing, summarization, coding,
and everyday information tasks. As AI-generated content becomes more common,
users need practical ways to evaluate whether a response is useful and
understand why it has a particular style, level of detail, or structure.
This is especially valuable when users have limited AI interactions available
and need to obtain useful results efficiently.

\section{Proposed Solution%
  \BvrHint{\ldots Software Proposition}}

\subsection{Overview}

Build a \textbf{web-based AI Response Analysis and Prompt Engineering
Platform} with the following components:

\begin{itemize}[leftmargin=0.7cm]
    \item \textbf{AI response analyzer:} evaluates relevance, clarity,
    conciseness, completeness, structure, usefulness, and adherence to
    requested constraints.
    \item \textbf{Response diagnosis:} identifies observable strengths and
    weaknesses such as excessive length, missing information, poor structure,
    unnecessary detail, or failure to follow requested format.
    \item \textbf{Response explanation:} provides a user-facing explanation
    of why the response has its observable characteristics by relating the
    response to prompt instructions and requested requirements.
    \item \textbf{Prompt analyzer:} identifies task, topic, audience, context,
    constraints, desired length, tone, and output format.
    \item \textbf{Prompt optimizer:} suggests changes to the prompt when the
    response does not satisfy the user's objective.
    \item \textbf{Comparison dashboard:} compares responses from different
    prompt versions and shows how their observable characteristics differ.
\end{itemize}

\subsection{Core workflow}

\begin{enumerate}[leftmargin=0.7cm]
    \item User enters a prompt and generates an AI response.
    \item The system analyzes the response for relevance, clarity,
    conciseness, completeness, structure, and adherence to the user's
    requirements.
    \item The dashboard identifies strengths, weaknesses, and possible
    response problems.
    \item The response explanation relates observable characteristics of the
    answer to identifiable instructions in the prompt.
    \item If the response is unsatisfactory, the system analyzes the prompt
    and suggests targeted improvements.
    \item The user can generate a revised response and compare it with the
    original response.
    \item The dashboard presents the comparison as a learning loop showing
    how prompt changes can affect observable response characteristics.
\end{enumerate}

\subsection{Operational constraints for an educational setting}

\begin{itemize}[leftmargin=0.7cm]
    \item Keep the project focused on text-based generative AI response
    analysis and user-facing explanation, with prompt engineering as a
    supporting improvement feature.
    \item Use an existing AI API instead of developing or training a large
    language model.
    \item Avoid claiming that individual words causally determine specific
    model outputs; explanations will be presented as likely effects based on
    prompt structure and observable response changes.
    \item Keep the first implementation small enough for a three-person team
    to complete within one semester.
    \item Prioritize software correctness, usability, measurable response
    analysis, response explanation, and practical experimentation over
    advanced AI research.
\end{itemize}

\section{Solution Approach%
  \BvrHint{\ldots Engineering focus}}

This is an engineering course project, so the focus will be on building a
practical software platform around an existing generative AI service. The
main contribution is the integration of response analysis, response
diagnosis, user-facing explanation, prompt analysis, and response comparison
into one application.

\subsection{Response analysis and prompt assistance%
  \BvrHint{\ldots non-research}}

Response analysis will be the primary analysis component of the system.
It will examine observable characteristics of generated answers, including:

\begin{itemize}[leftmargin=0.7cm]
    \item \textbf{Relevance:} whether the response addresses the user's
    actual request and avoids unnecessary information.
    \item \textbf{Clarity:} whether the response is understandable and uses
    an appropriate level of language for the requested audience.
    \item \textbf{Conciseness:} whether the response avoids unnecessary
    repetition and matches the requested level of detail.
    \item \textbf{Completeness:} whether the important parts of the requested
    task have been addressed.
    \item \textbf{Structure:} whether headings, bullets, steps, examples, or
    other requested formats are used appropriately.
    \item \textbf{Constraint adherence:} whether the response follows
    requirements such as word limits, tone, format, or audience.
    \item \textbf{Response diagnosis:} identifies observable problems and
    gives practical suggestions for improvement.
\end{itemize}

The response explanation will then connect these observations with the
prompt. For example, the system may identify that a response became detailed
because the prompt requested a detailed explanation, or that the response
was difficult for a beginner because no audience or complexity level was
specified. These are observable, user-facing explanations rather than claims
about hidden model reasoning.

Prompt assistance will identify task, topic, audience, context, constraints,
length, tone, and output format, and can suggest improvements when the
response analysis shows that the original request was insufficient.

\subsection{Web architecture%
  \BvrHint{\ldots web-based solution}}

A typical 3-tier web architecture:

\begin{itemize}[leftmargin=0.7cm]
    \item \textbf{Frontend:} responsive interface for entering prompts,
    viewing AI responses, displaying response-quality scores, showing
    response explanations and detected issues, and comparing responses.
    \item \textbf{Backend API:} handles AI API requests, response analysis,
    response diagnosis, response explanation, prompt analysis, prompt
    optimization, comparison logic, and history.
    \item \textbf{Database:} stores prompts, prompt versions, AI responses,
    selected objectives, analysis results, and timestamps.
\end{itemize}

\subsection{CI/CD and fast delivery enabler}

CI/CD will be used to:

\begin{itemize}[leftmargin=0.7cm]
    \item Automatically build the application and run tests on every change.
    \item Validate backend API and core prompt-processing functionality.
    \item Produce repeatable deployment artifacts.
    \item Support incremental development and safe integration of features.
\end{itemize}

\section{Evaluation Criterion%
  \BvrHint{\ldots measurable and attributable}}

We define success primarily using measurable criteria related to AI
response quality, response alignment, explanation usefulness, and the
ability of the platform to help users diagnose and improve unsatisfactory
responses.

\subsection{Primary evaluation metric}

\textbf{AI response quality and alignment:} measure whether the platform
can consistently identify important observable response characteristics and
whether generated responses satisfy the user's requested objective and
constraints.

\begin{itemize}[leftmargin=0.7cm]
    \item Evaluate response relevance, clarity, conciseness, completeness,
    structure, and compliance with requested constraints.
    \item Check whether identified response strengths and weaknesses are
    consistent with the actual response.
    \item Check whether the response explanation is supported by observable
    prompt and response characteristics.
    \item Compare original and revised responses to determine whether a
    targeted prompt change improved the selected response criteria.
\end{itemize}

\subsection{Secondary evaluation metrics}

\begin{itemize}[leftmargin=0.7cm]
    \item \textbf{Response quality:} relevance, clarity, conciseness,
    completeness, structure, and usefulness.
    \item \textbf{Response alignment:} degree to which the response follows
    requested length, format, audience, tone, and other constraints.
    \item \textbf{Response diagnosis accuracy:} whether detected strengths,
    weaknesses, and constraint mismatches correspond to the actual response.
    \item \textbf{Explanation usefulness:} whether users understand the
    observable reasons associated with the response characteristics.
    \item \textbf{Prompt impact:} whether targeted prompt changes produce
    measurable changes in response characteristics.
    \item \textbf{Efficiency:} reduction in unnecessary repeated prompts
    needed to obtain a satisfactory response during controlled tests.
\end{itemize}

\subsection{Pilot validation plan%
  \BvrHint{\ldots high-level}}

\begin{itemize}[leftmargin=0.7cm]
    \item Prepare common tasks such as explanation, summarization, coding,
    question answering, and summarizing text.
    \item Generate representative AI responses and evaluate them using the
    response-analysis dashboard.
    \item Test whether the system correctly identifies issues such as
    excessive length, missing information, poor structure, or unmet format
    requirements.
    \item Ask a small group of users whether the response explanations are
    understandable and whether the diagnosis helps them improve an answer.
    \item Test selected prompt changes and compare the resulting responses
    to determine whether the platform identifies meaningful improvements.
\end{itemize}

\section{Scalability%
  \BvrHint{\ldots with foresight}}

The platform should be scalable in both its software architecture and its
ability to support different AI response types, evaluation criteria, and
improvement workflows.

\begin{enumerate}[leftmargin=0.7cm]
    \item \textbf{Scalable (theoretically):} the system can support growth by:
    \begin{itemize}[leftmargin=0.7cm]
        \item separating frontend, backend, AI integration, and database
        components,
        \item using pagination for conversation and prompt history,
        \item indexing common database queries,
        \item allowing additional response metrics, explanation criteria,
        and prompt objectives to be added as modules,
        \item supporting multiple AI providers through a common backend
        interface in future versions.
    \end{itemize}

    \item \textbf{Operationally deployable (practically):} the system should
    run using:
    \begin{itemize}[leftmargin=0.7cm]
        \item standard web hosting,
        \item a managed or local database,
        \item containerized deployment where appropriate,
        \item CI/CD for repeatable builds and releases.
    \end{itemize}
\end{enumerate}

\section{Engine availability heuristic}

A mature set of tools and services is available to implement the platform:

\begin{itemize}[leftmargin=0.7cm]
    \item Web framework such as React or Next.js for the frontend.
    \item Python and FastAPI for backend REST APIs.
    \item A generative AI API for response generation.
    \item Natural-language processing and text-processing libraries for
    prompt and response analysis.
    \item PostgreSQL or another suitable database for prompt and response
    history.
    \item Visualization libraries for quality scores and prompt comparisons.
    \item Testing tools for backend, API, and prompt-analysis validation.
    \item Docker and CI/CD tools for deployment and repeatable development.
\end{itemize}

We will use mature components to focus on AI response analysis, response
explanation, prompt improvement, usability, and system integration rather than
developing a language model or AI infrastructure from scratch.

\section{Project scope and deliverables%
  \BvrHint{\ldots iteration-friendly}}

\subsection{Initial deliverable%
  \BvrHint{\ldots first iteration}}

\begin{itemize}[leftmargin=0.7cm]
    \item Prompt input and AI response generation.
    \item Core response-quality analysis for relevance, clarity,
    conciseness, completeness, structure, and constraint adherence.
    \item Response strengths, weaknesses, and response-diagnosis display.
    \item Basic response explanation connecting observable characteristics
    with prompt requirements.
    \item Basic prompt component identification.
    \item Initial prompt improvement suggestions.
    \item Database storage for prompts and responses.
    \item Basic frontend and backend integration.
    \item Automated tests for the main workflow.
\end{itemize}

\subsection{Subsequent deliverables}

\begin{itemize}[leftmargin=0.7cm]
    \item Detailed response explanation and response diagnosis.
    \item Prompt and response comparison with quality metrics.
    \item Response analysis for multiple task types such as explanation,
    coding, summarization, and question answering.
    \item Conversation and response history.
    \item Prompt evolution showing how successive changes affect responses.
    \item Prompt A/B testing for selected use cases.
    \item Prompt templates and learning suggestions.
    \item Optional support for multiple AI providers if time permits.
\end{itemize}

\section{Risks and mitigations}

\begin{itemize}[leftmargin=0.7cm]
    \item \textbf{Overclaiming AI explanations:} do not claim access to hidden
    model reasoning; explain only observable relationships between prompt
    structure and response characteristics.
    \item \textbf{Subjective response quality:} use clearly defined metrics
    and controlled comparisons rather than claiming a single universal
    quality score.
    \item \textbf{API limitations or cost:} use controlled test cases,
    configurable models, and minimize unnecessary API calls.
    \item \textbf{Project scope becoming too large:} prioritize response
    analysis, response explanation, response comparison, and basic prompt
    improvement before optional features.
    \item \textbf{Data privacy:} avoid storing unnecessary personal or
    sensitive conversation content and provide clear controls for stored
    history.
    \item \textbf{Model variability:} record the selected model and relevant
    request information when comparing responses so that results can be
    interpreted correctly.
    \item \textbf{Evaluation ambiguity:} define test prompts, objectives,
    response criteria, and comparison procedures before final evaluation.
\end{itemize}

\section{Summary}

This project proposes a deployable web platform centered on understanding
and improving AI-generated responses. The system evaluates responses for
relevance, clarity, conciseness, completeness, structure, usefulness, and
adherence to requested constraints. It then provides a user-facing diagnosis
and explanation of observable response characteristics.

Prompt engineering is used as a supporting mechanism: the system identifies
important prompt components and suggests targeted changes when the response
does not satisfy the user's goal. Users can compare original and revised
responses to see whether those changes improved the result. The overall
learning loop is:

\textbf{Prompt -> AI Response -> Response Analysis -> Explanation ->
Prompt Improvement -> Revised Response -> Comparison}

The project remains focused on software engineering rather than developing a
new language model. It combines frontend development, backend APIs, AI
integration, natural-language processing, database management, visualization,
testing, and deployment into a manageable three-person semester project.

\end{document}
